\chapter{MACHINE LEARNING \& GENERATIVE AI}

This block focuses on modern machine learning frameworks and generative models, with emphasis on deep learning architectures and training methodologies.

\begin{center}
\textbf{How to design, train, evaluate, and fine-tune advanced machine learning and generative models using state-of-the-art frameworks}
\end{center}

\vspace{0.5cm}

\section*{1. FRAMEWORKS}

\subsection*{1.1 PyTorch}

\begin{itemize}
\item Dynamic computational graphs
\item Pythonic design and flexibility
\item Strong support for research and experimentation
\end{itemize}

Example:

\begin{minted}{python}
import torch

x = torch.tensor([1.0, 2.0], requires_grad=True)
y = x**2
y.sum().backward()
\end{minted}

\subsection*{1.2 TensorFlow}

\begin{itemize}
\item Static and dynamic graph execution
\item Production-ready deployment
\item High-level API: Keras
\end{itemize}

\subsection*{1.3 JAX}

\begin{itemize}
\item Functional programming paradigm
\item Just-In-Time (JIT) compilation
\item Automatic differentiation and vectorization
\end{itemize}

Key transformation:

\[
\text{grad}(f(x)) = \frac{df}{dx}
\]

\vspace{0.5cm}

\section*{2. MODELS}

\subsection*{2.1 Transformers (LLMs)}

Transformers are based on the self-attention mechanism:

\[
\text{Attention}(Q,K,V) = \text{softmax}\left(\frac{QK^T}{\sqrt{d_k}}\right)V
\]

Key components:

\begin{itemize}
\item Self-attention
\item Multi-head attention
\item Positional encoding
\item Feedforward networks
\end{itemize}

Applications:

\begin{itemize}
\item Natural language processing
\item Code generation
\item Conversational AI
\end{itemize}

\vspace{0.3cm}

\subsection*{2.2 Variational Autoencoders (VAE)}

Generative models based on probabilistic encoding:

\[
z \sim q_\phi(z|x), \quad x \sim p_\theta(x|z)
\]

Loss function:

\[
\mathcal{L} = \mathbb{E}_{q_\phi(z|x)}[\log p_\theta(x|z)] - D_{KL}(q_\phi(z|x) \| p(z))
\]

Key ideas:

\begin{itemize}
\item Latent space representation
\item Probabilistic inference
\item Reconstruction + regularization
\end{itemize}

\vspace{0.3cm}

\subsection*{2.3 Diffusion Models}

Models that learn to reverse a noise process:

Forward process:

\[
q(x_t | x_{t-1})
\]

Reverse process:

\[
p_\theta(x_{t-1} | x_t)
\]

Training objective:

\[
\min_\theta \mathbb{E} \left[ \| \epsilon - \epsilon_\theta(x_t, t) \|^2 \right]
\]

Applications:

\begin{itemize}
\item Image generation
\item Audio synthesis
\item Scientific data generation
\end{itemize}

\vspace{0.5cm}

\section*{3. TECHNIQUES}

\subsection*{3.1 Fine-Tuning}

Adapting a pretrained model to a specific task:

\begin{itemize}
\item Transfer learning
\item Domain adaptation
\item Parameter-efficient methods (LoRA, adapters)
\end{itemize}

Update rule:

\[
\theta_{new} = \theta_{old} - \eta \nabla_\theta \mathcal{L}
\]

\vspace{0.3cm}

\subsection*{3.2 Model Evaluation}

Key metrics depend on the task:

\begin{itemize}
\item Classification: accuracy, precision, recall, F1-score
\item Generation: perplexity, BLEU, ROUGE
\item Image models: FID score
\end{itemize}

General loss minimization:

\[
\min_\theta \mathcal{L}(\theta)
\]

\vspace{0.5cm}

\section*{FINAL CONNECTION}

All components of this block connect through:

\begin{enumerate}
\item \textbf{Representation:} neural networks and latent spaces
\item \textbf{Learning:} optimization of parameters
\item \textbf{Generation:} sampling from learned distributions
\item \textbf{Evaluation:} measuring performance and generalization
\end{enumerate}

\vspace{0.5cm}

\section*{STUDY ROADMAP}

\subsection*{Phase 1 --- Frameworks}
\begin{itemize}
\item PyTorch fundamentals
\item TensorFlow/Keras basics
\item JAX transformations
\end{itemize}

\subsection*{Phase 2 --- Core Models}
\begin{itemize}
\item Transformers
\item VAEs
\item Diffusion models
\end{itemize}

\subsection*{Phase 3 --- Training Techniques}
\begin{itemize}
\item Optimization algorithms (SGD, Adam)
\item Fine-tuning strategies
\end{itemize}

\subsection*{Phase 4 --- Evaluation}
\begin{itemize}
\item Metrics and validation
\item Model comparison
\end{itemize}

\subsection*{Phase 5 --- Applications}
\begin{itemize}
\item NLP systems
\item Generative AI pipelines
\item Hybrid quantum-classical models
\end{itemize}